% GENERATED FILE. Edit canonical lessons, then run npm run generate:books.
% canonical-source-hash: fnv1a64:e59d6274147a93c2
% canonical-lessons: TA-C344-arakkan, TA-C344-tevatai, TA-C344-mantiravati, TA-C344-capi, TA-C344-urumaru

\chapter{Magic and Monsters}
\label{ch:ta-magic-and-monsters}
\hlchaptermodality{344}

\begin{chapteropening}
\textbf{By the end of this chapter:} \emph{I can tell a fairy tale with magic in it.}

Complete the last lesson of chapter 344: I can tell a fairy tale with magic in it.
\end{chapteropening}

\section[\texorpdfstring{\ta{அரக்கன்} (arakkaṉ)}{அரக்கன் (arakkaṉ)}]{\ta{அரக்கன்} (arakkaṉ) — a monster; a demon}

\label{lesson:TA-C344-arakkan}



\begin{quote}
Before the new one: say the Tamil for to make up, then the Tamil for a proverb.
\end{quote}

\subsection*{You'll want to know: \ta{அரக்கன்}}
\textbf{\ta{அரக்கன்}} — \emph{arakkaṉ} — \textquotedblleft{}a monster; a demon\textquotedblright{}.

\begin{quote}
\ta{கதையில்} \ta{ஒரு} \ta{அரக்கன்} \ta{காட்டில்} \ta{வாழ்ந்தான்}.
\end{quote}

\begin{quote}
\emph{kataiyil oru arakkaṉ kāṭṭil vāḻntāṉ.}
\end{quote}

\begin{quote}
In the story a monster lived in the forest.
\end{quote}

The female is \ta{அரக்கி}.

\begin{cousinweb}
\textbf{\ta{அரக்கன்}} (\emph{arakkaṉ}) is usually traced to Sanskrit \emph{rākṣasa}, \textquotedblleft{}demon\textquotedblright{} (CDIAL 10672), with the Tamil \ta{அ}- before a loan in r-.
\end{cousinweb}

\begin{grammarlens}[title={What you've built}]
One more word for telling what happened.
\end{grammarlens}

\subsection*{Your turn}
\begin{itemize}
\raggedright
  \item \emph{Say it:} \emph{arakkaṉ}
  \item \emph{Say it:} \emph{arakkaṉ}, once more
  \item \emph{Say it:} \emph{kataiyil oru arakkaṉ kāṭṭil vāḻntāṉ.} — In the story a monster lived in the forest.
  \item \emph{Recall:} say the Tamil for to make up, then the Tamil for a proverb, then say \emph{arakkaṉ} again
\end{itemize}

\begin{tcolorbox}[breakable,colback=teal!4,colframe=teal!35!black,title={Before you move on}]
What does \ta{அரக்கன்} mean? (\textquotedblleft{}A monster; a demon\textquotedblright{}.) Say it once more.
\end{tcolorbox}
\section[\texorpdfstring{\ta{தேவதை} (tēvatai)}{தேவதை (tēvatai)}]{\ta{தேவதை} (tēvatai) — a fairy; an angel}

\label{lesson:TA-C344-tevatai}



\begin{quote}
Before the new one: say the Tamil for a proverb, then the Tamil for a monster.
\end{quote}

\subsection*{You'll want to know: \ta{தேவதை}}
\textbf{\ta{தேவதை}} — \emph{tēvatai} — \textquotedblleft{}a fairy; an angel\textquotedblright{}.

\begin{quote}
\ta{ஒரு} \ta{தேவதை} \ta{ஏழைச்} \ta{சிறுமிக்கு} \ta{உதவினாள்}.
\end{quote}

\begin{quote}
\emph{oru tēvatai ēḻaic ciṟumikku utaviṉāḷ.}
\end{quote}

\begin{quote}
A fairy helped the poor girl.
\end{quote}

Also of a very kind person.

\begin{cousinweb}
\textbf{\ta{தேவதை}} (\emph{tēvatai}) is borrowed from Sanskrit \emph{dēvatā}, \textquotedblleft{}divinity\textquotedblright{} (CDIAL 6530).
\end{cousinweb}

\begin{grammarlens}[title={What you've built}]
One more word for telling what happened.
\end{grammarlens}

\subsection*{Your turn}
\begin{itemize}
\raggedright
  \item \emph{Say it:} \emph{tēvatai}
  \item \emph{Say it:} \emph{tēvatai}, once more
  \item \emph{Say it:} \emph{oru tēvatai ēḻaic ciṟumikku utaviṉāḷ.} — A fairy helped the poor girl.
  \item \emph{Recall:} say the Tamil for a proverb, then the Tamil for a monster, then say \emph{tēvatai} again
\end{itemize}

\begin{tcolorbox}[breakable,colback=teal!4,colframe=teal!35!black,title={Before you move on}]
What does \ta{தேவதை} mean? (\textquotedblleft{}A fairy; an angel\textquotedblright{}.) Say it once more.
\end{tcolorbox}
\section[\texorpdfstring{\ta{மந்திரவாதி} (mantiravāti)}{மந்திரவாதி (mantiravāti)}]{\ta{மந்திரவாதி} (mantiravāti) — a sorcerer; a magician}

\label{lesson:TA-C344-mantiravati}



\begin{quote}
Before the new one: say the Tamil for a monster, then the Tamil for a fairy.
\end{quote}

\subsection*{You'll want to know: \ta{மந்திரவாதி}}
\textbf{\ta{மந்திரவாதி}} — \emph{mantiravāti} — \textquotedblleft{}a sorcerer; a magician\textquotedblright{}.

\begin{quote}
\ta{மந்திரவாதி} \ta{இளவரசனை} \ta{ஒரு} \ta{தவளையாக} \ta{மாற்றினான்}.
\end{quote}

\begin{quote}
\emph{mantiravāti iḷavaracaṉai oru tavaḷaiyāka māṟṟiṉāṉ.}
\end{quote}

\begin{quote}
The sorcerer turned the prince into a frog.
\end{quote}

From \ta{மந்திரம்}, a spell.

\begin{cousinweb}
\textbf{\ta{மந்திரவாதி}} (\emph{mantiravāti}) is borrowed from Sanskrit \emph{mantravādin}, \textquotedblleft{}a reciter of spells, an enchanter\textquotedblright{} (Monier-Williams).
\end{cousinweb}

\begin{grammarlens}[title={What you've built}]
One more word for telling what happened.
\end{grammarlens}

\subsection*{Your turn}
\begin{itemize}
\raggedright
  \item \emph{Say it:} \emph{mantiravāti}
  \item \emph{Say it:} \emph{mantiravāti}, once more
  \item \emph{Say it:} \emph{mantiravāti iḷavaracaṉai oru tavaḷaiyāka māṟṟiṉāṉ.} — The sorcerer turned the prince into a frog.
  \item \emph{Recall:} say the Tamil for a monster, then the Tamil for a fairy, then say \emph{mantiravāti} again
\end{itemize}

\begin{tcolorbox}[breakable,colback=teal!4,colframe=teal!35!black,title={Before you move on}]
What does \ta{மந்திரவாதி} mean? (\textquotedblleft{}A sorcerer; a magician\textquotedblright{}.) Say it once more.
\end{tcolorbox}
\section[\texorpdfstring{\ta{சபி} (capi)}{சபி (capi)}]{\ta{சபி} (capi) — to curse}

\label{lesson:TA-C344-capi}



\begin{quote}
Before the new one: say the Tamil for a fairy, then the Tamil for a sorcerer.
\end{quote}

\subsection*{You'll want to know: \ta{சபி}}
\textbf{\ta{சபி}} — \emph{capi} — \textquotedblleft{}to curse\textquotedblright{}.

\begin{quote}
\ta{கோபமான} \ta{முனிவர்} \ta{மன்னரைச்} \ta{சபித்தார்}.
\end{quote}

\begin{quote}
\emph{kōpamāṉa muṉivar maṉṉaraic capittār.}
\end{quote}

\begin{quote}
The angry sage cursed the king.
\end{quote}

Past (I …): \ta{சபித்தேன்}, \emph{capittēṉ}.

The noun is \ta{சாபம்}, a curse.

\begin{cousinweb}
\textbf{\ta{சபி}} (\emph{capi}) is a Tamil verb built on Sanskrit \emph{śapati}, \textquotedblleft{}curses, swears\textquotedblright{} (CDIAL 12289).
\end{cousinweb}

\begin{grammarlens}[title={What you've built}]
One more word for telling what happened.
\end{grammarlens}

\subsection*{Your turn}
\begin{itemize}
\raggedright
  \item \emph{Say it:} \emph{capi}
  \item \emph{Say it:} \emph{capi}, once more
  \item \emph{Say it:} \emph{kōpamāṉa muṉivar maṉṉaraic capittār.} — The angry sage cursed the king.
  \item \emph{Recall:} say the Tamil for a fairy, then the Tamil for a sorcerer, then say \emph{capi} again
\end{itemize}

\begin{tcolorbox}[breakable,colback=teal!4,colframe=teal!35!black,title={Before you move on}]
What does \ta{சபி} mean? (\textquotedblleft{}To curse\textquotedblright{}.) Say it once more.
\end{tcolorbox}
\section[\texorpdfstring{\ta{உருமாறு} (urumāṟu)}{உருமாறு (urumāṟu)}]{\ta{உருமாறு} (urumāṟu) — to change shape; to transform}

\label{lesson:TA-C344-urumaru}



\begin{quote}
Before the new one: say the Tamil for a sorcerer, then the Tamil for to curse.
\end{quote}

\subsection*{You'll want to know: \ta{உருமாறு}}
\textbf{\ta{உருமாறு}} — \emph{urumāṟu} — \textquotedblleft{}to change shape; to transform\textquotedblright{}.

\begin{quote}
\ta{அந்தத்} \ta{தவளை} \ta{ஒரு} \ta{இளவரசனாக} \ta{உருமாறியது}.
\end{quote}

\begin{quote}
\emph{antat tavaḷai oru iḷavaracaṉāka urumāṟiyatu.}
\end{quote}

\begin{quote}
That frog changed into a prince.
\end{quote}

Past (I …): \ta{உருமாறினேன்}, \emph{urumāṟiṉēṉ}.

What it becomes takes -\ta{ஆக}.

\begin{cousinweb}
\textbf{\ta{உருமாறு}} (\emph{urumāṟu}) is a Tamil compound built from \ta{உரு}, \textquotedblleft{}form, shape\textquotedblright{} (DEDR 657, which asks whether it comes from Sanskrit \emph{rūpa}), and \ta{மாறு}, \textquotedblleft{}to become changed\textquotedblright{} (DEDR 4834).
\end{cousinweb}

\begin{grammarlens}[title={What you've built}]
One more word for telling what happened.
\end{grammarlens}

\subsection*{Your turn}
\begin{itemize}
\raggedright
  \item \emph{Say it:} \emph{urumāṟu}
  \item \emph{Say it:} \emph{urumāṟu}, once more
  \item \emph{Say it:} \emph{antat tavaḷai oru iḷavaracaṉāka urumāṟiyatu.} — That frog changed into a prince.
  \item \emph{Recall:} say the Tamil for a sorcerer, then the Tamil for to curse, then say \emph{urumāṟu} again
\end{itemize}

\begin{tcolorbox}[breakable,colback=teal!4,colframe=teal!35!black,title={Before you move on}]
What does \ta{உருமாறு} mean? (\textquotedblleft{}To change shape; to transform\textquotedblright{}.) Say it once more.
\end{tcolorbox}
